\chapter{Resultados}
\label{chap:resultados}

\section{Distancia de fibra}

La Figura~\ref{fig:exp1_distance} muestra el barrido de distancia. La probabilidad de detección analítica cae casi exponencialmente con la distancia, como se espera para una fibra con atenuación expresada en dB/km. El QBER se mantiene bajo durante la mayor parte del barrido y aumenta abruptamente en el punto de mayor distancia, donde supera el umbral de referencia de \(11\,\%\). Este comportamiento indica que, bajo los parámetros elegidos, el enlace conserva margen operativo hasta distancias del orden de varias decenas de kilómetros, pero se vuelve inestable cerca del extremo superior del barrido.

\begin{figure}[H]
  \centering
  \includegraphics[width=\textwidth]{exp1_distance_sweep.png}
  \caption{QBER y probabilidad de detección en función de la distancia de fibra para BB84 \textit{time-bin}.}
  \label{fig:exp1_distance}
\end{figure}

La Figura~\ref{fig:exp1_skr} presenta la tasa de llave secreta estimada. La tasa cae gradualmente con la distancia mientras el QBER permanece por debajo del umbral, y se anula en la práctica cuando el error excede el margen de seguridad del modelo. La escala logarítmica permite observar tanto la degradación progresiva como la caída final.

\begin{figure}[H]
  \centering
  \includegraphics[width=0.92\textwidth]{exp1_skr_distance.png}
  \caption{Tasa de llave secreta estimada en función de la distancia.}
  \label{fig:exp1_skr}
\end{figure}

\section{Sensibilidad del detector}

La Figura~\ref{fig:exp2_detector} resume el impacto de eficiencia y conteos oscuros a \SI{50}{\kilo\meter}. En el barrido de eficiencia, el QBER permanece por debajo del umbral para todos los puntos simulados. La SKR cambia con la eficiencia de detector, aunque el resultado muestra variaciones punto a punto asociadas al carácter estocástico y al tamaño finito de muestra. Esto refuerza que el barrido sirve como indicador de tendencia y sensibilidad, no como predicción determinista punto por punto.

En el barrido de conteos oscuros, el QBER también se mantiene bajo en las condiciones simuladas. La degradación esperada por conteos oscuros no domina fuertemente este conjunto de parámetros, lo que sugiere que, para el banco de pruebas inicial, la eficiencia y la acumulación de eventos válidos pueden ser restricciones más visibles que los conteos oscuros dentro del rango elegido. En una implementación física, este resultado debería revisarse con ventanas temporales, jitter y tasas de repetición reales.

\begin{figure}[H]
  \centering
  \includegraphics[width=\textwidth]{exp2_detector_sensitivity.png}
  \caption{QBER y SKR frente a eficiencia de detector y conteos oscuros a \SI{50}{\kilo\meter}.}
  \label{fig:exp2_detector}
\end{figure}

\section{Visibilidad interferométrica}

La Figura~\ref{fig:exp3_visibility} muestra el barrido de visibilidad. A medida que la visibilidad mejora, el QBER tiende a disminuir y la SKR se ubica en valores más altos. El efecto no es perfectamente monótono porque la simulación usa llaves cortas y un número limitado de repeticiones; aun así, la relación física principal se conserva: la calidad interferométrica es una variable crítica para una codificación \textit{time-bin}.

Este resultado es importante para el diseño del banco de pruebas. La estabilidad de fase, el balance de caminos ópticos y la resolución temporal del detector no son detalles secundarios, sino condiciones necesarias para que la base de superposición produzca resultados confiables. En consecuencia, el interferómetro y la sincronización deberían considerarse subsistemas de alta prioridad experimental.

\begin{figure}[H]
  \centering
  \includegraphics[width=0.92\textwidth]{exp3_visibility.png}
  \caption{QBER y SKR en función de la visibilidad del interferómetro a \SI{30}{\kilo\meter}.}
  \label{fig:exp3_visibility}
\end{figure}

\section{Estados señuelo}

La Figura~\ref{fig:exp4_decoy} compara la estimación sin estados señuelo con la cota analítica decoy. Ambas curvas disminuyen con la distancia. La curva decoy queda por debajo de la estimación simple porque aplica una cota más conservadora sobre la contribución segura de eventos monofotónicos. Por lo tanto, el resultado no implica que los estados señuelo empeoren el sistema; indica que la tasa segura certificable bajo ese modelo es menor que una estimación simple que no contempla la misma exposición a pulsos multifotónicos.

Para un banco de pruebas real, la lectura correcta es que los estados señuelo agregan complejidad experimental y de procesamiento, pero son relevantes para acercar la estimación de tasa a un modelo de seguridad más realista con fuentes coherentes débiles. El paso siguiente sería implementar selección aleatoria de intensidades y estimación estadística dentro del flujo protocolar, en lugar de calcular la cota como postproceso analítico.

\begin{figure}[H]
  \centering
  \includegraphics[width=0.92\textwidth]{exp4_decoy_impact.png}
  \caption{Comparación de SKR modelo sin estados señuelo y cota analítica con estados señuelo.}
  \label{fig:exp4_decoy}
\end{figure}

\section{Síntesis de resultados}

\begin{table}[H]
  \centering
  \begin{tabular}{p{0.25\textwidth}p{0.64\textwidth}}
    \toprule
    Barrido & Conclusión de diseño \\
    \midrule
    Distancia & La atenuación reduce la probabilidad de detección y termina anulando la SKR cuando el QBER supera el umbral. \\
    Detector & La eficiencia afecta la tasa útil; los conteos oscuros del rango probado no dominan el error bajo estos parámetros. \\
    Visibilidad & La calidad interferométrica es crítica para \textit{time-bin}; visibilidad baja aumenta QBER y reduce SKR. \\
    Estados señuelo & La cota decoy es más conservadora y representa mejor el riesgo de pulsos multifotónicos en fuentes coherentes débiles. \\
    \bottomrule
  \end{tabular}
  \caption{Resumen de conclusiones experimentales para el diseño del banco de pruebas.}
  \label{tab:sintesis_resultados}
\end{table}
